\begin{abstract}
Soft-tissue imbalance after total knee arthroplasty can compromise joint stability, and surgeons therefore balance the soft tissues during the operation.
Patient-specific knee models could support these intraoperative decisions by predicting how a proposed adjustment changes joint restraint.
Such predictions require ligament parameters, which must be inferred from measurements that are available during surgery.
We therefore asked which intraoperative measurements constrain these parameters.
We built a rigid-body model of a cruciate-retaining knee replacement with nonlinear ligament bundles and simulated joint gaps and anterior--posterior, varus--valgus, and internal--external laxity across flexion.
We then used the local Fisher information of the ligament parameters to compare measurement sets.
In addition, we used bounded nonlinear least squares to evaluate parameter inference and predictions under held-out load cases.
In our model, joint gaps alone primarily constrained medial collateral ligament reference strain and provided little information about stiffness.
Adding multiplanar laxity measurements across flexion improved parameter identifiability, although stiffness and reference strain remained coupled.
Calibration to the full examination improved predictions of held-out passive laxity, but parameter errors remained comparable to those of the prior center.
Using the true ligament attachment sites reduced the normalized parameter root-mean-square error from 0.31 to 0.17.
These findings demonstrate that accurate laxity predictions do not necessarily imply accurate ligament properties and highlight the importance of both measurement selection and attachment-site accuracy in patient-specific knee model calibration.
\end{abstract}

\keywords{total knee arthroplasty; soft-tissue balancing; knee laxity; model calibration; identifiability}
\linenumbers
\modulolinenumbers[1]
%% ------------------------------------------------------------------------------
\section{Introduction}

In 2020, about 595 million people worldwide had osteoarthritis, most commonly of the knee, and knee cases are projected to rise by about \SI{75}{\percent} by 2050~\autocite{gbd2023osteoarthritis}.
Total knee arthroplasty is the standard treatment for end-stage disease.
Yet a systematic review estimated that about \SI{10}{\percent} of patients remain dissatisfied after the operation~\autocite{defrance2023dissatisfied}.
Instability is a common reason for revision, including within two years of surgery~\autocite{brown2022evolving,athwal2014clinical}.
To provide stability while allowing knee motion, surgeons balance the soft tissues during the operation.
They assess joint gaps and laxity through flexion.
To correct excessive tightness or laxity, they may then adjust bone resections, component position, or insert thickness, or release tight soft tissues.
Gap-balancing techniques use a tensioned soft-tissue envelope to guide component positioning, whereas measured-resection techniques initially rely on anatomical landmarks~\autocite{sheth2017surgical,daines2014gap}.
However, joint gaps describe only the opening under distraction.
They do not fully describe the resistance of the knee to anterior--posterior (AP), varus--valgus (VV), and internal--external (IE) loading~\autocite{markolf1976stiffness,athwal2014clinical}.

Computational knee models can quantify the mechanical consequences of surgical choices such as component rotation~\autocite{thompson2011biomechanical}, and model studies show that both component alignment and ligament properties influence tibiofemoral contact forces~\autocite{smith2016influence}.
Patient-specific models may therefore help surgeons evaluate balancing adjustments.
Sensor-instrumented trial components have already been used to quantify compartment loading during surgery~\autocite{gustke2014increased}.
A model-calibration workflow would additionally require the applied loads and the resulting motion during an instrumented examination.
A model calibrated to these measurements could then predict passive laxity at other flexion angles.
After further development and validation, it could also, for instance, estimate how a thicker insert would alter restraint before the surgeon selects that insert.
In such models, ligament stiffness and reference strain govern the force response, but neither is directly measurable during surgery.
Their inference is therefore a means of personalizing model predictions rather than a clinical endpoint in itself~\autocite{ewing2016estimating}.

Computational models have been calibrated to measured knee laxity to estimate specimen-specific ligament parameters~\autocite{harris2016combined,ewing2016estimating}.
However, close agreement with measured laxity does not establish that the fitted parameters are unique or accurate.
Different parameter combinations may produce similar laxity~\autocite{gutenkunst2007universally}.
In addition, uncertainty in the attachment sites may be absorbed by the fitted parameters, because these sites can be difficult to determine from MRI~\autocite{rachmat2014generating}.
Local identifiability analysis provides a framework to quantify which parameter directions measurements constrain under specified model and noise assumptions~\autocite{raue2009structural,gutenkunst2007universally}.
To our knowledge, however, this framework has not been applied to intraoperative laxity measurements in knee arthroplasty.
Assessing parameter identifiability is an important step before interpreting a calibrated model or using it to guide an adjustment.

In this simulation study, we used a cruciate-retaining knee replacement model to assess how joint gaps and multiplanar laxity measurements across flexion constrain ligament stiffness and reference strain.
We compared measurement sets using local Fisher information, and we used bounded nonlinear least squares to evaluate parameter inference and passive laxity predictions under held-out load cases (\autoref{fig:overview}).
We also examined how attachment mismatch affects the interpretation of the fitted parameters.
Our contribution is therefore a formal identifiability assessment of intraoperative measurements for knee model personalization, together with a demonstration that agreement with measured laxity alone does not establish tissue properties.

\begin{figure}[H]
\centering
\includegraphics{fig1_overview.pdf}
\caption{\textbf{A virtual laxity examination links ligament properties to measurable joint motion.}
(A)~Anterior, medial, posterior, and lateral views of the right-knee model in the reference pose at full extension.
Toroidal condyles of the femoral component articulate with a dished tibial insert, and nine nonlinear bundles represent the medial collateral ligament (MCL), lateral collateral ligament (LCL), posterior cruciate ligament (PCL), and capsule (CAP).
In the anterior view, the femoral origin $O_\mathrm{F}$ lies midway between the condyles on the flexion axis, and the tibial origin $O_\mathrm{T}$ lies midway between the lowest points of the insert dishes.
Filled and open circles mark the nominal femoral and tibial attachments of \autoref{tab:bundles}, and each view labels the bundles nearest the viewer. Bundle prefixes denote anterior (a), intermediate (i), anterolateral (al), posteromedial (pm), and posterolateral (pl). Bones are omitted.
The four loading conditions were (1) joint distraction for the medial and lateral gaps and (2) AP, (3) VV, and (4) IE laxity tests, each at \SIlist{0;30;60;90}{\degree} flexion.
(B)~The current study generated synthetic knees, simulated their examinations, and assessed parameter inference by Fisher information and least squares.
In the intended clinical use, an instrumented examination would update the model, which would then predict the effect of an adjustment.}
\label{fig:overview}
\end{figure}

%% ------------------------------------------------------------------------------
\section{Methods}\label{Methods}

\subsection{Rigid-body knee model}

We modeled the femur and tibia of a right cruciate-retaining knee replacement as rigid bodies, following established multibody knee models~\autocite{wismans1980three,blankevoort1991articular,bei2004multibody}.
The femoral frame had its origin midway between the condyles on the flexion axis, with $x$ anterior, $y$ medial, and $z$ proximal (\autoref{fig:overview}A).
The tibial frame had the same axes in the reference pose at full extension, and its origin lay midway between the lowest points of the two insert dishes, \SI{25}{\milli\meter} below the femoral origin.
For a prescribed flexion angle $\alpha$, a tibial point $\ten{p}_\mathrm{T}$ in the tibial frame was located in the femoral frame at
\begin{equation}
\ten{p}_\mathrm{F} = \ten{R}_y(\alpha)\left[\ten{R}_x(\beta)\,\ten{R}_z(\gamma)\,\ten{p}_\mathrm{T} + \ten{u}\right].
\label{eq:pose}
\end{equation}
Here, $\ten{u} = [u_\mathrm{AP}, u_\mathrm{ML}, u_\mathrm{PD}]$ is the vector from the femoral to the tibial origin in the flexed femoral frame, $\beta$ is varus rotation about the floating axis perpendicular to the flexion axis and the tibial long axis, and $\gamma$ is internal rotation about the tibial long axis.
This rotation sequence matches the joint coordinate system of \textcite{grood1983joint}, whereas $\ten{u}$ is resolved in the flexed femoral frame rather than along its translation axes.
At each flexion angle, static equilibrium determined these five coordinates, $\ten{q} = [\ten{u}, \beta, \gamma]$.

The implant geometry is listed in \autoref{tab:implant}.
The femoral component comprised two single-radius toroidal condyles, and the insert comprised two paraboloidal dishes, with a more conforming medial dish.
Each grid point $i$ of the insert, with area $A$, carried the normal contact force of an elastic foundation~\autocite{blankevoort1991articular,bei2004multibody},
\begin{equation}
\ten{f}_i = k_\mathrm{c}\,A\,\langle -\delta_i \rangle_+\,\ten{n}_i.
\label{eq:contact}
\end{equation}
Here, $\delta_i$ is the signed separation from the toroidal condyle, $\ten{n}_i$ is the outward normal of the condyle, $\langle x\rangle_+$ is the positive part, smoothed over \SI{0.01}{\milli\meter}, and $k_\mathrm{c}$ is the foundation modulus.

\begin{table}[H]
\centering
\caption{Implant geometry and contact parameters. Medial--lateral positions are measured from the femoral and tibial origins. The foundation modulus is close to the value of about \SI{270}{\newton\per\cubic\milli\meter} given by the elastic-foundation formula for an \SI{8}{\milli\meter} polyethylene layer with the material properties used by \textcite{bei2004multibody}.}
\label{tab:implant}
\begin{tabular}{llc}
\toprule
Component & Quantity & Value\\
\midrule
Femoral component & sagittal radius of each condyle & \SI{25}{\milli\meter}\\
 & coronal radius of each condyle & \SI{20}{\milli\meter}\\
 & medial--lateral position of the condyle centers & $\pm$\SI{22}{\milli\meter}\\
Tibial insert & AP radius of the medial and lateral dishes & \SIlist{45;80}{\milli\meter}\\
 & medial--lateral radius of both dishes & \SI{26}{\milli\meter}\\
 & articular surface, AP by medial--lateral & \SI{50}{\milli\meter} by \SI{71}{\milli\meter}\\
Contact & grid spacing on each dish & \SI{1.5}{\milli\meter}\\
 & foundation modulus $k_\mathrm{c}$ & \SI{250}{\newton\per\cubic\milli\meter}\\
\bottomrule
\end{tabular}
\end{table}


Nine tension-only bundles followed the force--strain law of \textcite{blankevoort1991ligament},
\begin{equation}
f(\varepsilon) =
\begin{cases}
0, & \varepsilon \le 0,\\[2pt]
\dfrac{k\,\varepsilon^2}{4\varepsilon_\mathrm{l}}, & 0 < \varepsilon < 2\varepsilon_\mathrm{l},\\[6pt]
k\,(\varepsilon - \varepsilon_\mathrm{l}), & \varepsilon \ge 2\varepsilon_\mathrm{l},
\end{cases}
\qquad
\varepsilon = \frac{l}{l_0} - 1,
\qquad
l_0 = \frac{l_\mathrm{r}}{1 + \varepsilon_\mathrm{r}}.
\label{eq:ligament}
\end{equation}
Here, $k$ is the bundle stiffness in force per unit strain, $\varepsilon_\mathrm{l} = 0.03$ defines the toe region, $l$ is the current length, and $l_\mathrm{r}$ is the length in the reference pose.
The reference strain $\varepsilon_\mathrm{r}$ is the bundle strain in the reference pose, so that $k$ and $\varepsilon_\mathrm{r}$ together determine the reference-pose force.
The attachment sites followed anatomical studies~\autocite{laprade2007anatomy,laprade2003posterolateral,claes2013anatomy,amis2006anatomy} and are listed in \autoref{tab:bundles}.
Because $l_\mathrm{r}$ was computed from the attachment sites of each knee, perturbed attachment sites (\autoref{sec:population}) changed $l_\mathrm{r}$ and $l_0$, and $\varepsilon_\mathrm{r}$ always denoted the reference-pose strain of that knee.
The bundles formed four groups, and the bundles of each group shared one stiffness and one reference strain (\autoref{tab:bundles}).
This yields eight unknown parameters, $\ten{\vartheta} = [k_\mathrm{MCL}; \varepsilon_{\mathrm{r},\mathrm{MCL}}; k_\mathrm{LCL}; \varepsilon_{\mathrm{r},\mathrm{LCL}}; k_\mathrm{PCL}; \varepsilon_{\mathrm{r},\mathrm{PCL}}; k_\mathrm{CAP}; \varepsilon_{\mathrm{r},\mathrm{CAP}}]$.

\begin{table}[H]
\centering
\caption{Ligament bundles and their nominal attachment sites. Coordinates are in millimeters, with $x$ anterior, $y$ medial, and $z$ proximal. Femoral sites are given in the femoral frame and tibial sites in the tibial frame, whose origin lies on the joint line \SI{25}{\milli\meter} below the femoral origin at full extension. The nominal sites are based on anatomical studies~\autocite{laprade2007anatomy,laprade2003posterolateral,claes2013anatomy,amis2006anatomy}. The reference length $l_\mathrm{r}$ is the bundle length in the reference pose.}
\label{tab:bundles}
\small
\begin{tabular}{llccr}
\toprule
Group & Bundle & Femoral site & Tibial site & $l_\mathrm{r}$\\
 & & $(x, y, z)$ [\si{\milli\meter}] & $(x, y, z)$ [\si{\milli\meter}] & [\si{\milli\meter}]\\
\midrule
MCL & anterior & ($4$, $40$, $4$) & ($12$, $32$, $-60$) & 90\\
MCL & intermediate & ($0$, $41$, $4$) & ($2$, $33$, $-57$) & 86\\
MCL & posterior & ($-6$, $40$, $8$) & ($-20$, $33$, $-20$) & 55\\
LCL & single & ($-3$, $-42$, $4$) & ($-24$, $-38$, $-38$) & 70\\
PCL & anterolateral & ($4$, $9$, $-3$) & ($-24$, $2$, $-11$) & 44\\
PCL & posteromedial & ($-4$, $13$, $0$) & ($-27$, $5$, $-15$) & 47\\
Capsule & posteromedial & ($-28$, $24$, $14$) & ($-27$, $24$, $-7$) & 46\\
Capsule & posterolateral & ($-28$, $-24$, $14$) & ($-27$, $-24$, $-7$) & 46\\
Capsule & anterolateral & ($-5$, $-41$, $6$) & ($-2$, $-35$, $-6$) & 38\\
\bottomrule
\end{tabular}
\end{table}


Static equilibrium required the generalized forces of contact, ligaments, and external loads to vanish,
\begin{equation}
\ten{Q}(\ten{q}) = \ten{Q}_\mathrm{contact}(\ten{q}) + \ten{Q}_\mathrm{ligament}(\ten{q}) + \ten{Q}_\mathrm{external}(\ten{q}) = \ten{0}.
\label{eq:equilibrium}
\end{equation}
Each component of $\ten{Q}$ is the generalized force conjugate to one coordinate of $\ten{q}$.
Using the kinematics in \autoref{eq:pose}, we projected the resultant force onto the axes of $\ten{u}$ and its moment about the tibial origin onto the rotation axes of $\beta$ and $\gamma$, while the prescribed flexion supplied the reaction about the flexion axis.
We solved \autoref{eq:equilibrium} with a Newton method.
The Jacobian was computed by forward differences, and each Newton step was limited to \SI{2}{\milli\meter} for AP and medial--lateral translation, \SI{1}{\milli\meter} for proximal translation, and \SIlist{0.05;0.1}{\radian} for varus and internal rotation.
Iterations stopped when the largest residual fell below \SI{1e-4}{\newton}, with moments divided by a length scale of \SI{25}{\milli\meter}.
If a solve did not reach a residual below \SI{0.01}{\newton}, the load was applied in 3, 8, and then 20 equal increments.
A check of path dependence is reported in the Supplementary Material.
We implemented the model and all analyses in Python with NumPy and SciPy~\autocite{harris2020array,virtanen2020scipy}.

\subsection{Virtual laxity examination}

Each laxity test applied a \SI{50}{\newton} axial compression that seated the femur in the insert, with external loads expressed in the tibial frame.
At flexion angles of \SIlist{0;30;60;90}{\degree}, we applied anterior and posterior forces of \SI{100}{\newton} at \SI{5}{\milli\meter} below the joint line, varus and valgus moments of \SI{10}{\newton\meter}, and internal and external moments of \SI{3}{\newton\meter}.
Laxity was the change of the corresponding coordinate from the compression-only equilibrium at the same flexion angle.
For the gap test, we replaced the compression by a distraction force of \SI{100}{\newton} at the center of each compartment.
The AP, medial--lateral, and IE coordinates were held at their compression-only values, while proximal translation and varus rotation were free.
The medial and lateral joint gaps were the minimum separations between the condyle and the insert in each compartment.
The full examination thus comprised 32 measurements, namely six laxities (anterior, posterior, varus, valgus, internal, and external) and two gaps at each of the four flexion angles.

\subsection{Virtual knee population and measurement noise}\label{sec:population}

We sampled the stiffnesses log-uniformly and the reference strains uniformly within the ranges of \autoref{tab:prior}.
Collateral reference strains ranged from zero to positive values, and PCL and capsule reference strains from negative to positive values.
These ranges were model assumptions rather than measured patient distributions.
To represent uncertainty in locating the attachment sites~\autocite{rachmat2014generating}, we perturbed each attachment coordinate by Gaussian noise with an assumed standard deviation of \SI{1.5}{\milli\meter}.
To characterize the forward model, we simulated \num{10000} knees (Supplementary \autoref{fig:forward}).
Independently, we simulated 500 test knees, and the examinations of 499 of them converged.
We used all converged test knees for the Fisher analysis and the first 50 converged knees for least squares.
Measurement noise was Gaussian with assumed standard deviations of \SI{0.5}{\milli\meter} for translations and gaps and \SI{0.5}{\degree} for rotations.
Each test knee received one random draw of this noise, and the same noisy measurements were used in all least-squares fits.

\begin{table}[H]
\centering
\caption{Prior ranges of the ligament parameters. Stiffness is per bundle and per unit strain. Stiffness was sampled log-uniformly and reference strain uniformly.}
\label{tab:prior}
\begin{tabular}{lcccc}
\toprule
 & MCL & LCL & PCL & Capsule \\
\midrule
Bundles & 3 & 1 & 2 & 3\\
Stiffness $k$ [\si{\kilo\newton}] & 1.5--5.5 & 1.0--4.0 & 2.5--9.0 & 0.5--3.0\\
Reference strain $\varepsilon_\mathrm{r}$ [\si{\percent}] & 0 to 8 & 0 to 8 & $-4$ to 4 & $-4$ to 4\\
\bottomrule
\end{tabular}
\end{table}

\subsection{Local Fisher information of measurement sets}

We compared nine measurement sets, from gaps at \SIlist{0;90}{\degree} only to the full examination.
For each test knee, we computed the sensitivity matrix $\ten{J} = \partial \ten{y}/\partial \hat{\ten{\vartheta}}$ of the measurements $\ten{y}$ with respect to the normalized parameters $\hat{\ten{\vartheta}}$.
The normalization mapped the prior range of each parameter to the unit interval, logarithmically for stiffness and linearly for reference strain.
We computed $\ten{J}$ by central differences with a step of $0.02$, kept inside the prior.
The uniform prior on the unit interval has variance $1/12$, so a Gaussian with the same variance has precision 12.
With this Gaussian approximation of the prior, the linearized posterior covariance is
\begin{equation}
\ten{C} = \left(\ten{J}^\mathrm{T}\ten{\Sigma}^{-1}\ten{J} + 12\,\ten{I}\right)^{-1},
\label{eq:fisher}
\end{equation}
where $\ten{\Sigma}$ is the diagonal noise covariance and $\ten{J}^\mathrm{T}\ten{\Sigma}^{-1}\ten{J}$ is the Fisher information~\autocite{raue2009structural,gutenkunst2007universally}.
We quantified the local reduction in parameter standard deviation relative to this moment-matched Gaussian prior using the identifiability index $I_j = 1 - \sqrt{C_{jj}}/\sqrt{1/12}$.
It is zero when the measurements add nothing to the prior and one when they determine the parameter exactly.
From $\ten{C}$, we also obtained the local correlation between log stiffness and reference strain of each ligament group.
The attachment sites of each knee were held fixed, so our analysis is conditional on known attachment sites and contains no term for attachment uncertainty.

\subsection{Calibration by bounded nonlinear least squares}

We calibrated the model to the noisy examination of each of the 50 knees.
The fit error was half the sum of squared noise-normalized residuals between the simulated and noisy measurements.
We minimized it with the trust-region reflective algorithm~\autocite{branch1999subspace}, the SciPy method for bound-constrained nonlinear least squares~\autocite{virtanen2020scipy}.
We fitted the model with the nominal attachment sites, so the perturbed attachment sites of each knee acted as an unknown attachment mismatch.
Parameters were optimized in normalized coordinates within the unit box, with forward-difference Jacobians (relative step 0.01) and the default SciPy tolerances.
Each fit started from the prior center and from two points drawn uniformly from $[0.1, 0.9]^8$, with a budget of 500 residual evaluations excluding those for the Jacobians.
For each knee, we retained the fit with the lowest final fit error.
We called two starts of a knee distinct when their fitted parameters differed by more than 0.2 in at least one normalized parameter and both reached a fit error within one unit of the lowest fit error of that knee.
To test whether laxity testing improves calibration beyond gaps, we repeated the fits using only the four gaps at \SIlist{0;90}{\degree}, with the same knees, noise, bounds, and starts.
To isolate the effect of attachment mismatch, we also repeated the full-examination fits with the true attachment sites of each knee.

\subsection{Predictions under held-out load cases}

We used the fitted parameters to predict laxity at \SIlist{15;45;75}{\degree} flexion and under three combined loads at \SI{45}{\degree}, none of which entered calibration.
The combined loads were a varus moment of \SI{10}{\newton\meter} with an internal moment of \SI{3}{\newton\meter}, a valgus moment with an external moment of the same magnitudes, and an anterior force of \SI{100}{\newton} with an internal moment of \SI{3}{\newton\meter}, each with the \SI{50}{\newton} compression.
The reference values used the true parameters and true attachment sites of each knee without measurement noise.
Predictions used the same attachment sites as the corresponding fit.
Parameter error was the root-mean-square error (RMSE) of the normalized parameters, computed with log-scaled prior bounds for stiffness and linear bounds for reference strain and pooled over knees and, for the overall value, over all eight parameters.
Held-out AP, VV, and IE errors were RMSEs pooled with equal weight over knees and nine outputs each.
These outputs were the two load directions at three flexion angles and the corresponding coordinate under each of the three combined loads.
As baselines, we predicted with the prior center and with the true parameters combined with the nominal attachment sites.
The latter isolates the effect of attachment mismatch and is not a zero-error reference.

%% ------------------------------------------------------------------------------
\section{Results}

\subsection{Contact-grid refinement had little effect on the nominal-knee examination}

Refining the contact grid from \SI{1.5}{\milli\meter} to \SI{0.5}{\milli\meter} changed no measurement of the nominal knee by more than \SI{0.04}{\milli\meter} or \SI{0.04}{\degree}, which is more than ten times less than the measurement noise.
Equilibrium converged for all held-out load cases of the 50 knees used for least squares.

\subsection{Multiplanar testing strengthens parameter constraints across flexion}

Joint gaps alone provided local, prior-regularized constraints mainly on the MCL reference strain (\autoref{fig:identifiability}A).
With gaps at \SIlist{0;90}{\degree}, as in conventional gap balancing~\autocite{sheth2017surgical}, the identifiability index of every stiffness stayed below 0.1, whereas that of the MCL reference strain was 0.39.
Each added laxity test then strengthened the constraints on specific ligament groups.
Adding AP testing to gaps and VV testing increased the PCL stiffness index from 0.26 to 0.54.
Adding IE testing instead increased the LCL stiffness index from 0.12 to 0.39 and the capsule stiffness index from 0.29 to 0.62.
The full examination at extension alone gave a PCL stiffness index of only 0.02, compared with 0.64 at four flexion angles.

\begin{figure}[H]
\centering
\includegraphics{fig2_identifiability.pdf}
\caption{\textbf{Multiplanar testing strengthens local parameter constraints, while stiffness and reference strain remain coupled.}
(A)~Median local identifiability index of each ligament parameter for nine measurement sets across 499 virtual knees with known attachment sites.
The index is $1 - \sigma_\mathrm{post}/\sigma_\mathrm{prior}$ (\autoref{eq:fisher}), so zero denotes no reduction of the prior standard deviation and one an exactly determined parameter.
Gaps alone mainly constrained the MCL reference strain, whereas adding VV, AP, and IE laxity tests across flexion raised the indices of most parameters.
$k$ denotes bundle stiffness and $\varepsilon_\mathrm{r}$ reference strain of the medial collateral (MCL), lateral collateral (LCL), and posterior cruciate ligament (PCL) and the capsule (CAP).
Gaps are medial and lateral openings, and AP, VV, and IE denote anterior--posterior, varus--valgus, and internal--external laxity.
``Full exam'' includes gaps and all three laxity tests, and ``0'' and ``0/90'' restrict the flexion angles to \SI{0}{\degree} and to \SIlist{0;90}{\degree}, whereas the other rows use \SIlist{0;30;60;90}{\degree}.
(B)~Median local correlation between log stiffness and reference strain of each ligament group, computed from the covariance in normalized parameter coordinates.
The MCL and PCL correlations became more strongly negative as laxity tests were added.}
\label{fig:identifiability}
\end{figure}

With the full examination, the collateral reference-strain indices reached 0.74 (MCL) and 0.70 (LCL), compared with stiffness indices of 0.48 and 0.43.
In contrast, the capsule indices were similar for stiffness (0.64) and reference strain (0.61).
Log stiffness and reference strain remained negatively correlated in the local covariance approximation, with correlations of $-0.84$ for the MCL and $-0.82$ for the PCL (\autoref{fig:identifiability}B).

\subsection{Agreement with measured laxity does not establish accurate ligament parameters}

With the full examination and the nominal attachment sites, least squares fitted the measured laxity closely but did not infer the true parameters accurately (\autoref{fig:calibration}A).
All 150 fits stopped before the evaluation budget, when the change in fit error or in the parameter step fell below the default tolerances, after at most 68 residual evaluations.
In all 50 knees, the fitted parameters gave a lower fit error than the true parameters with the nominal attachment sites, by a median factor of 5.0, to a median of 29.7.
The overall normalized parameter RMSE was nevertheless 0.31, similar to the 0.29 of the prior center.
The collateral reference strains improved on the prior center, whereas the errors of the collateral stiffnesses exceeded those of the prior center.
In 23 of the 50 knees, the three starts ended more than 0.2 apart in at least one normalized parameter, and in 6 of these knees the starts were distinct by our fit-error criterion.
Agreement among three starts does not establish a unique global optimum, because the multistart comparison detects, but cannot exhaustively rule out, alternative fits.

\begin{figure}[H]
\centering
\includegraphics{fig3_calibration.pdf}
\caption{\textbf{Multiplanar calibration improves held-out laxity prediction, while accurate parameter inference depends on the attachment sites.}
Results are for 50 test knees with three bounded least-squares starts per knee.
(A)~RMSE of the normalized parameters $\hat\vartheta$, that is, as a fraction of the log-scaled stiffness range or the linear reference-strain range of \autoref{tab:prior}.
Fits to the full examination with nominal attachment sites had errors similar to those of the prior center, whereas fits with the true attachment sites had the smallest errors.
Fits to gaps at \SIlist{0;90}{\degree} alone had the largest errors.
(B,C)~RMSE of predicted AP translation and of VV and IE rotation at \SIlist{15;45;75}{\degree} flexion and under three combined loads at \SI{45}{\degree}, none of which entered calibration, compared with noise-free simulations using the true parameters and attachment sites.
Full-examination fits predicted these load cases better than the true parameters combined with the nominal attachment sites.}
\label{fig:calibration}
\end{figure}

The four-gap fits had a median fit error of 0.17, yet their normalized parameter RMSE was 0.40, which exceeded that of the prior center.
In 43 of the 50 knees, the starts were distinct by the same criterion.
With only four measurements for eight parameters, the gap fits could not locally constrain all parameters without prior information and did not identify unique parameters.

With the true attachment sites, the full-examination fits had a median fit error of 12.3, close to the value of 12 expected for 32 measurements and eight fitted parameters under a local linear approximation.
Their normalized parameter RMSE fell from 0.31 to 0.17, with collateral reference-strain errors of 0.13 and 0.10.
The starts were distinct in only 3 of the 50 knees.

Despite their parameter error, the fits with nominal attachment sites predicted held-out passive laxity well (\autoref{fig:calibration}B,C).
Compared with the prior center, they reduced the held-out AP error from \SI{1.9}{\milli\meter} to \SI{0.8}{\milli\meter}, the VV error from \SI{1.4}{\degree} to \SI{0.7}{\degree}, and the IE error from \SI{4.8}{\degree} to \SI{1.7}{\degree}.
These errors were also smaller than those of the true parameters with the nominal attachment sites.
In pooled RMSE, fits to gaps alone predicted only slightly better than the prior center, and fits with the true attachment sites predicted best (per-knee variability in \autoref{tab:heldout_iqr}).
Calibration with the nominal attachment sites took a median of \SI{5.2}{\minute} per knee for all three starts on one core.

%% ------------------------------------------------------------------------------
\section{Discussion}

In our study, joint gaps, the measurement on which gap-balancing techniques rely~\autocite{sheth2017surgical}, poorly constrained ligament stiffness.
Knees with different collateral stiffnesses could therefore present similar gaps yet respond differently to varus, anterior, and rotational loads.
Multiplanar testing across flexion strengthened the constraints on most parameters.
This result is consistent with the plane- and flexion-dependent contributions of ligaments and capsule documented in native-knee sectioning studies~\autocite{markolf1976stiffness}.
These comparisons support testing across planes and flexion angles, but they do not establish a minimal examination.
Because attachment sites were held fixed in the Fisher analysis, these indices describe information conditional on known attachment sites, whereas the least-squares fits show how attachment mismatch degrades parameter inference.

The persistent coupling between stiffness and reference strain indicates that our examination weakly constrained some parameter combinations.
Under the force--strain law of \textcite{blankevoort1991ligament} and with a single load magnitude per test, a higher stiffness combined with a lower reference strain can produce ligament forces, and hence laxities, similar to those of a lower stiffness with a higher reference strain.
Such coupled directions are typical of nonlinear models with many parameters~\autocite{gutenkunst2007universally,raue2009structural,transtrum2015perspective}.
They imply that laxity-based calibrations~\autocite{harris2016combined,ewing2016estimating} may return one of several parameter sets that fit similarly well.
Parameter uncertainty should therefore accompany model calibration.

Fits with the nominal attachment sites predicted held-out synthetic laxity well, although their parameter error was no smaller than that of the prior center.
Supplying the true attachment sites improved the inferred parameters, which is consistent with compensation for attachment mismatch in the fits with nominal attachment sites.
Separately, multiple parameter sets fitted the gaps closely.
Thus, agreement with measured laxity alone established neither unique nor accurate parameters.
Coefficients that absorb attachment mismatch should be treated as effective model parameters rather than direct tissue measurements.

Our results also suggest how inference could be improved, although we did not test these remedies.
Measuring laxity in all three planes is not sufficient by itself, because the full examination already included gaps and all three laxity tests and the coupling persisted.
Additional load magnitudes may help separate stiffness from reference strain by sampling a broader range of ligament recruitment.
Locating the attachment sites more precisely, for example from imaging or navigation, could reduce the attachment mismatch that biased these fits.
Where neither is possible, Bayesian calibration with informative priors can report the remaining uncertainty rather than a single parameter set.

The intended clinical use is to personalize a model during an instrumented examination with trial components so that it can help evaluate the mechanical consequences of a proposed adjustment.
With further validation, such a workflow may help compare changes in component position, insert thickness, or soft-tissue release, as computational studies have related soft-tissue modifications to multiplanar laxity in other implant designs~\autocite{amiri2012computational}.
Moreover, parameters that compensate for attachment mismatch may reproduce the current examination yet give unreliable predictions after an intervention changes the load path.
Our calibration took several minutes per knee, so real-time surgical use remains untested.
Simulation-trained regression has enabled rapid personalization of implanted hip-capsule models, which suggests a possible route to faster inference, but we did not evaluate that approach here~\autocite{ananthakrishnan2024instantaneous}.

Our study has a few limitations.
The model is a rigid-body simplification with straight, tension-only ligament bundles, no wrapping, and a simplified posterolateral restraint, which may explain its asymmetric rotational laxity.
Adding model detail would not by itself resolve the identifiability problem, because it may introduce further unknowns~\autocite{transtrum2015perspective}.
The four ligament groups shared parameters across their bundles, whereas real bundles differ in stiffness and recruitment~\autocite{amis2006anatomy}.
The geometry was generic, and only attachment sites were varied, although bone morphology and component alignment also differ between patients~\autocite{smith2016influence}.
All data were synthetic, and the model was not compared with measured laxity of replaced knees.
The preload and independent measurement-noise levels were assumed rather than measured intraoperatively.
Doubling the noise lowered the indices but preserved the ranking of measurement sets (Supplementary Material).
Our results therefore concern the passive mechanical calibration of a model, not active postoperative function or patient outcomes, and the responses to releases or insert changes were not evaluated.
The necessary next step is testing against cadaveric laxity measurements of the kind used for specimen-specific calibration~\autocite{harris2016combined,ewing2016estimating}, ideally combined with tensile tests of the same ligaments~\autocite{butler1986comparison}.

%% ------------------------------------------------------------------------------
\section{Conclusion}

We formalized the question of which intraoperative measurements constrain ligament properties with a local identifiability analysis and least-squares calibration of a knee replacement model.
Our study shows that joint gaps alone poorly constrained ligament stiffness, whereas multiplanar laxity testing across flexion strengthened local parameter constraints.
Agreement with measured laxity alone did not establish unique or accurate ligament parameters, and supplying the true attachment sites reduced parameter error.
Our findings support broader measurement collection when personalizing knee models for future intraoperative decision support.
Clinical validation and faster, reliable calibration are needed before such models can guide real-time balancing decisions.

